\chapter{Introduction}
\begin{center}
\large \textbf{Evolution of Key Establishment Paradigms: From Public-Key Encryption to Key Encapsulation Mechanism}
\end{center}
\vspace{1em}

Since the birth of public-key cryptography, its theoretical abstraction and engineering practice have continuously evolved through mutual promotion and compromise. Before exploring why modern secure protocols have comprehensively embraced the Key Encapsulation Mechanism (KEM), it is essential to review the semantics of classic key establishment primitives and the bottlenecks they encountered in real-world deployments.

\section{Limitations of Classic Primitives: Diffie-Hellman and PKE}

Early public-key cryptography sought to address two independent yet deeply intertwined challenges: establishing a shared secret over an insecure channel, and enabling non-interactive confidential communication. In 1976, Diffie and Hellman proposed their landmark key exchange protocol (DH) based on the discrete logarithm problem~\cite{diffie2022new}. The elegance of DH lies in its use of homomorphic mathematical structures (e.g., $g^{xy} = (g^x)^y$), allowing communicating parties to ``aggregate'' a shared secret out of public interactions. Furthermore, when instantiated with ephemeral keys, DH naturally provides \textit{forward secrecy}. However, its fundamental limitation is the strict requirement for interactive communication.

In parallel, general-purpose Public-Key Encryption (PKE), epitomized by RSA, offered an alternative paradigm. PKE functions as a digital ``envelope''---the sender uses the receiver's public key to encrypt an arbitrary message, which the receiver then decrypts. Under traditional semantic definitions ($c \gets \text{Enc}(pk, m)$), a PKE scheme enables non-interactive encryption but typically lacks native forward secrecy unless augmented with protocol-level key rotations. 

Moreover, these two paradigms diverge significantly in their security goals. DH is typically analyzed in the random oracle or generic group model, targeting Indistinguishability under Chosen Plaintext Attack (IND-CPA) style security for the resulting shared secret. Conversely, robust PKE schemes are designed to achieve Indistinguishability under Chosen Ciphertext Attack (IND-CCA) for arbitrary message spaces. In real-world network protocols, using PKE to directly process large volumes of application data is computationally prohibitive. Consequently, early engineering practices (such as RSA key transport in early TLS) relegated PKE to a fallback role: encrypting a short ``pre-master secret'' before switching to symmetric encryption. While functional, the powerful syntax of PKE---supporting arbitrary plaintext input---was largely wasted and often became a fertile ground for side-channel and chosen-ciphertext attacks.

\section{The Protocol-Centric Shift: KEM Abstraction and Engineering Synergy}

To bridge the gap between theoretical primitives and actual protocol requirements, the cryptographic community began to redefine the boundaries of key establishment. The genuine requirement of modern secure channels is remarkably specific: protocols do not need to encrypt meaningful text using public-key algorithms; they merely need to synchronize a short, high-entropy random bitstring between two endpoints.

Motivated by this observation, the concept of KEM was formalized, notably through Shoup's formalization of the KEM/DEM (Data Encapsulation Mechanism) composition paradigm in 2001~\cite{shoup2001proposal}. In the KEM syntax, the encapsulation algorithm $(K, c) \gets \text{Encaps}(pk)$ strips the caller of the ability to input a specific plaintext. Instead, the algorithm internally generates a random shared secret $K$ and outputs its encapsulation $c$. 

This subtle semantic shift triggered profound architectural improvements in modern protocol design. Primarily, it provided exceptional modularity underpinned by formal security guarantees. The KEM/DEM composition theorem states that combining an IND-CCA secure KEM with an IND-CCA secure DEM straightforwardly yields an IND-CCA secure hybrid encryption scheme. In modern protocols like the Hybrid Public Key Encryption standard (HPKE, RFC 9180)~\cite{rfc9180}, the KEM is solely responsible for outputting high-entropy secrets, the Key Derivation Function (KDF) securely binds this secret to communication transcripts and identities, and an Authenticated Encryption with Associated Data (AEAD) cipher manages the bulk data. This decoupling not only simplifies security proofs but also facilitates highly flexible cipher suite negotiations.

More importantly, KEM serves as the vital bridge to the post-quantum cryptography (PQC) era. Under the threat of quantum computing, classical DH and PKE algorithms will be rendered insecure. However, mainstream post-quantum public-key schemes---particularly lattice-based (e.g., LWE or NTRU), code-based, and isogeny-based constructions---often involve the deliberate introduction of mathematical noise. Forcing these noisy structures to perfectly encrypt a user-specified plaintext requires significant parameter redundancy and complex Error Correction Codes (ECC). In contrast, negotiating a noisy random value and aligning it via hashing into a shared key is highly natural and efficient.

Because the KEM syntax perfectly circumvents the requirement to ``carry exact plaintexts,'' it resonates flawlessly with the internal mechanics of PQC primitives. This explains why the National Institute of Standards and Technology (NIST), in its newly released post-quantum standard FIPS 203 (ML-KEM)~\cite{nist2024fips203}, completely abandoned the traditional PKE interface, standardizing KEM as the sole asymmetric paradigm. The transition from PKE to KEM is not merely a sign of cryptography maturing from theory to engineering; it is the fundamental architecture enabling a seamless migration to the post-quantum future.

\section{Scope and Assumptions}

This review assumes the reader's familiarity with foundational topics typically covered in a modern cryptography course. Specifically, we presume a working knowledge of semantic security definitions for public-key encryption, the Learning With Errors (LWE) and Short Integer Solution (SIS) problems underlying lattice-based cryptography.

Consequently, this paper does not re-introduce these underlying mathematical primitives from scratch. Instead, our scope is strictly focused on the KEM abstraction itself: its formal security properties, its transformational constructions (such as the Fujisaki-Okamoto transform for achieving CCA security), and its structural role in large-scale protocol deployment. Mathematical preliminaries are kept minimal, introducing only the requisite notations necessary to analyze KEM constructions and rigorously evaluate their security arguments.
